\section{The genus of a curve}
\label{section-genus}

\noindent
If $X$ is a smooth projective geometrically irreducible curve over a field $k$,
then we've previously defined the genus of $X$ as the dimension of
$H^1(X, \mathcal{O}_X)$, see
Picard Schemes of Curves, Definition \ref{pic-definition-genus}.
Observe that $H^0(X, \mathcal{O}_X) = k$ in this case, see
Varieties, Lemma \ref{varieties-lemma-regular-functions-proper-variety}.
Let us generalize this as follows.

\begin{definition}
\label{definition-genus}
Let $k$ be a field. Let $X$ be a proper scheme over $k$ having
dimension $1$ and $H^0(X, \mathcal{O}_X) = k$.
Then the {\it genus} of $X$ is $g = \dim_k H^1(X, \mathcal{O}_X)$.
\end{definition}

\noindent
This is sometimes called the {\it arithmetic genus} of $X$.
In the literature the arithmetic genus of a proper curve $X$
over $k$ is sometimes defined as
$$
p_a(X) = 1 - \chi(X, \mathcal{O}_X) =
1 - \dim_k H^0(X, \mathcal{O}_X) + \dim_k H^1(X, \mathcal{O}_X)
$$
This agrees with our definition when it applies because we assume
$H^0(X, \mathcal{O}_X) = k$. But note that
\begin{enumerate}
\item $p_a(X)$ can be negative, and
\item $p_a(X)$ depends on the base field $k$ and should be written $p_a(X/k)$.
\end{enumerate}
For example if $k = \mathbf{Q}$
and $X = \mathbf{P}^1_{\mathbf{Q}(i)}$ then
$p_a(X/\mathbf{Q}) = -1$ and $p_a(X/\mathbf{Q}(i)) = 0$.

\medskip\noindent
The assumption that $H^0(X, \mathcal{O}_X) = k$ in our definition has
two consequences. On the one hand, it means there is no confusion about
the base field. On the other hand, it implies the scheme $X$ is
Cohen-Macaulay and equidimensional of dimension $1$
(Lemma \ref{lemma-automatic}). If $\omega_X$ denotes the dualizing
module as in Lemmas \ref{lemma-duality-dim-1} and
\ref{lemma-duality-dim-1-CM} we see that
\begin{equation}
\label{equation-genus}
g = \dim_k H^1(X, \mathcal{O}_X) = \dim_k H^0(X, \omega_X)
\end{equation}
by duality, see Remark \ref{remark-rework-duality-locally-free}.

\medskip\noindent
If $X$ is proper over $k$ of dimension $\leq 1$ and $H^0(X, \mathcal{O}_X)$
is not equal to the ground field $k$, instead of using the arithmetic genus
$p_a(X)$ given by the displayed formula above we shall use the invariant
$\chi(X, \mathcal{O}_X)$. In fact, it is advocated in
\cite[page 276]{FAC} and \cite[Introduction]{Hirzebruch}
that we should call $\chi(X, \mathcal{O}_X)$ the arithmetic genus.

\begin{lemma}
\label{lemma-genus-base-change}
Let $k'/k$ be a field extension. Let $X$ be a proper scheme over $k$ having
dimension $1$ and $H^0(X, \mathcal{O}_X) = k$. Then $X_{k'}$ is a
proper scheme over $k'$
having dimension $1$ and $H^0(X_{k'}, \mathcal{O}_{X_{k'}}) = k'$.
Moreover the genus of $X_{k'}$ is equal to the genus of $X$.
\end{lemma}

\begin{proof}
The dimension of $X_{k'}$ is $1$ for example by
Morphisms, Lemma \ref{morphisms-lemma-dimension-fibre-after-base-change}.
The morphism $X_{k'} \to \Spec(k')$ is proper by
Morphisms, Lemma \ref{morphisms-lemma-base-change-proper}.
The equality $H^0(X_{k'}, \mathcal{O}_{X_{k'}}) = k'$ follows from
Cohomology of Schemes, Lemma
\ref{coherent-lemma-flat-base-change-cohomology}.
The equality of the genus follows from the same lemma.
\end{proof}

\begin{lemma}
\label{lemma-genus-gorenstein}
Let $k$ be a field. Let $X$ be a proper scheme over $k$ having
dimension $1$ and $H^0(X, \mathcal{O}_X) = k$. If $X$ is Gorenstein,
then
$$
\deg(\omega_X) = 2g - 2
$$
where $g$ is the genus of $X$ and $\omega_X$ is as in
Lemma \ref{lemma-duality-dim-1}.
\end{lemma}

\begin{proof}
Immediate from Lemma \ref{lemma-rr}.
\end{proof}

\begin{lemma}
\label{lemma-genus-smooth}
Let $X$ be a smooth proper curve over a field $k$
with $H^0(X, \mathcal{O}_X) = k$. Then
$$
\dim_k H^0(X, \Omega_{X/k}) = g
\quad\text{and}\quad
\deg(\Omega_{X/k}) = 2g - 2
$$
where $g$ is the genus of $X$.
\end{lemma}

\begin{proof}
By Lemma \ref{lemma-duality-dim-1} we have $\Omega_{X/k} = \omega_X$.
Hence the formulas hold by (\ref{equation-genus}) and
Lemma \ref{lemma-genus-gorenstein}.
\end{proof}






